\section*{الفصل \ref{ch:g1:living-or-not} --- حيٌّ أم غير حي؟}

\begin{solution}{exo:g1:living-or-not:1}
يولد، ويتغذى، وينمو، ويمكن أن يكون له صغار، ويموت. واليوم نرى
الهريرة تتغذى (تشرب الحليب وتأكل) وتنمو --- فهي تكبر أسبوعًا بعد
أسبوع.
\end{solution}

\begin{solution}{exo:g1:living-or-not:2}
(أ) حية --- خرجت من بيضة، وتتغذى على الرحيق، وقد تكون لها يرقات
صغيرة. (ب) غير حية --- صُنعت في مصنع، ولا تتغذى ولا تنمو أبدًا.
(ج) حية --- نمت من بذرة. (د) غير حية --- السحابة تكبر وتصغر،
لكنها لم تولد قط، ولا تتغذى، وليست لها صغار.
\end{solution}

\begin{solution}{exo:g1:living-or-not:3}
ليست حية الآن. لكنها كانت حية: فقد نمت على الشجرة، وغذّتها
الشجرة. وهي تأتي من \omterm{def:g1:living-or-not:living}{كائن حي} --- الشجرة.
\end{solution}

\begin{solution}{exo:g1:living-or-not:4}
اسأل أسئلة الحياة مجتمعة: أوُلد النهر؟ أيتغذى؟ أينمو؟ أيمكن أن
تكون له أنهار صغيرة؟ لا --- فالحركة والضجيج ليستا علامتي حياة
بذاتهما. النهر ليس حيًّا، وإن عاشت \emph{فيه} كائنات حية كثيرة.
\end{solution}

\begin{solution}{exo:g1:living-or-not:5}
حي: العنكبوت، والتفاحة على الشجرة. كان حيًّا (أو يأتي من \omterm{def:g1:living-or-not:living}{كائن
حي}): الجورب الصوفي (صوف من خروف)، وكوز الصنوبر الساقط (نما على
صنوبرة). \omterm{def:g1:living-or-not:nonliving}{لم يكن حيًّا قط}: قطرة المطر.
\end{solution}

\begin{solution}{exo:g1:living-or-not:6}
لا. اللهب يسقط في أسئلة الحياة: فقد أُشعل ولم يولد، والأهم أنه
ليست له ألسنة لهب صغيرة تكبر --- يمكن نقل اللهب إلى شمعة أخرى،
لكن ليست له صغار كالحلزون أو زهرة الأقحوان. (وهو سؤال وجيه: فاللهب
يخدع فعلًا علامتين من العلامات!)
\end{solution}

\begin{solution}{exo:g1:living-or-not:7}
مثلًا: كرسي خشبي (من شجرة)، وقميص قطني (من نبات القطن)، وبطانية
صوفية (من خروف)، وحذاء جلدي (من بقرة)، وعسل على الرف (من النحل).
\end{solution}

\begin{solution}{exo:g1:living-or-not:8}
الفولة تتغير: فتنتفخ، ثم يخرج منها جذير أبيض صغير، ثم برعم أخضر.
أما الحصاة فتبقى كما هي تمامًا. وهذا يدل على أن الفولة كانت حية
--- كائنًا حيًّا هادئًا ينتظر --- وأن الحصاة ليست حية.
\end{solution}

\begin{solution}{exo:g1:living-or-not:9}
أوُلد؟ الحلزون خرج من بيضة، والإنسان الآلي صُنع. أيمكن أن تكون له
صغار؟ قد تكون هناك حلزونات صغيرة، ولا يوجد آليون صغار. (والتغذي
يصلح أيضًا: الخس ينمي الحلزون، والبطارية لا تنمي الإنسان الآلي.)
\end{solution}

\begin{solution}{exo:g1:living-or-not:10}
الفطر يولد (يظهر من بدايات صغيرة في التربة)، ويتغذى (يأخذ غذاءه
من التربة ومن الأوراق الميتة)، وينمو --- في ليلة واحدة أحيانًا ---
ويصنع فطورًا جديدة، ويموت. فكل علامات الحياة حاضرة: إنه \omterm{def:g1:living-or-not:living}{كائن حي}،
ولو بلا أرجل وبلا أوراق خضراء.
\end{solution}
